\chapter{Physikalische Grundlagen}\label{chapter:physikalische_grundlagen}

Im folgenden Kapitel sollen die Grundlagen erl?utert werden, auf denen diese Vorlage und sp?tere Arbeit basiert.
Zuerst wird jedoch ein einziges Paper zitiert, um auch ein Literaturverzeichnis zu erstellen \cite{Oltmann_2010} and \cite{Oltmann_2012}.

\section{Erstes Unterkapitel}\label{sec:erstes_unterkapitel}

Hier steht das erste Unterkapitel des Kapitels \ref{chapter:physikalische_grundlagen}. Es tr?gt auch hier die Nummer \ref{sec:erstes_unterkapitel}.


\subsection{Erstes Unterunterkapitel}\label{subsec:erstes_unterunterkapitel}
Hier werden zwei Gleichungen dargestellt wobei die Gleichung (\ref{eq:kugelwelle}) die Kugelwelle beschreibt.

\begin{equation}
	\label{eq:kugelwelle}
		u_K(\mathbf{r},t)=\frac{u_0}{r} \cdot e^{ikr-\omega t},
\end{equation}

\begin{equation}
	\label{eq:ebene_welle}
		u_E(\mathbf{r},t)=u_0 \cdot e^{i\mathbf{kr}-\omega t}.
\end{equation}

Jetzt folgt eine EPS-Grafik durch den Befehl (\textbackslash PictureCaptionHeight\{...\}\{...\}\{...\}\{...\}). Dieser Befehl braucht vier ?bergabewerte: Den Dateinahmen, den Titel der Abbildung im Abbildungsverzeichnis, der Titel direkt unterhalb der Abbildung im Text, sowie die Bildh?he in beliebigen Einheiten, hier beispielsweise 35\,mm. Den Befehl gibt es analog f?r die Bildbreite (\textbackslash PictureCaptionWidth\{...\}\{...\}\{...\}\{...\}), wobei hier der vierte Parameter die Bildbreite ist. Der Befehl geht nur f?r EPS-Dateien, die Endung '.eps' im Dateinamen entf?llt. Die Grafik liegt im Bildordner, dieser wird im cls-file aktuel auf  den Ordner 'Images' definiert (\textbackslash def \textbackslash scr@bildverzeichnis\{images\}). Das Label der Grafik ist automatisch der Bildname. Da die Grafik zu gro? ist, springt sie auf die n?chste Seite.

\PictureCaptionHeight{beugung_einfachspalt}{Beugung einer ebenen Welle am Einfachspalt}{Beugung einer senkrecht auf einen Spalt treffenden ebenen Welle und Beispiel f?r ein resultierendes Beugungsbild.}{35mm}




\subsection{Zweites Unterunterkapitel}\label{sucsec:zweites_unterunterkapitel}

Jetzt kommen zwei Grafiken nebeneinander mit gleicher Bildh?he durch den Befehl (\textbackslash TwoPicturesHeight\{Name1\}\{Name2\}\{Titel Abb.Verz.\}\{Bild-unterschrift\}\{Abstand hor.\}\{Bildh?he\}) mit sechs ?bergabewerten: Zwei Bildnamen, die Unterschriften f?r Abbildungsverzeichnis und Bildunterschrift, sowie den horizontalen Abstand und die jeweilige Breite der Bilder.

\TwoPicturesHeight{airy_scheibchen}{rechtecksblende}{Beugungsbild Airy und Rechtecksblende}{Hier sieht man sehr deutlich zwei Bilder nebeneinander und den Unterschied zwischen Blendenformen}{10mm}{35mm}


Jetzt kommen drei Grafiken nebeneinander mit gleicher H?he, der Befehl funktioniert analog, das Bild erscheint aus Platzgr?nden eine Seite weiter.\newpage

\ThreePicturesHeight{beugung_einfachspalt_rotiert}{airy_scheibchen}{rechtecksblende}{Beugungsbild Einfachspalt Rotiert, Airy und Rechtecksblende}{Hier sieht man sehr deutlich drei Bilder nebeneinander und den Unterschied zwischen Blendenformen sowie den gedrehten Einfachspalt}{10mm}{25mm}

Und jetzt noch eine 2x2 Darstellung:

\FourPicturesBoxHeight{rechtecksblende}{airy_scheibchen}{beugung_einfachspalt_rotiert}{beugung_einfachspalt}{Vier Bilder als Block}{Hier sieht man sch?n vier Bilder neben- und untereinander im Block.}{5mm}{25mm}

Analoge Befehle gibt es f?r die jeweiligen Bildbreiten, dazu cls-file durchsehen. Weitere Definitionen sind gerne selbst zu erg?nzen.

Jetzt kommt eine beispielhafte  Tabelle. Hier muss jeder selbst anpassen um die gew?nschte Zeilen und Spaltenanzahl zu erhalten.

\begin{table}
\begin{center}
\caption[Kurzer Titel f?r Tabellenverzeichnis]{Titel f?r Tabelle selbst, welcher nun auch etwas l?nger sein darf. Bei mehreren Zeilen wird es unsch?n.}
\label{tab:tab1}
\begin{tabular}{ccc}
%\hline
\toprule[1.2pt]
\textbf{Bezeichnung 1} & \textbf{Bezeichnung 2} & \textbf{Bezeichnung 3} \\  %\hline
\midrule[1.2pt]\midrule[1.0pt]
Wert 11       & Wert 12     & Wert 13  	    \\
Wert 21       & Wert 22     & Wert 23  	    \\
Wert 31       & Wert 32     & Wert 33  	    \\
\hline
\end{tabular}
\end{center}
\end{table}

Hier kommt jetzt massenhaft Text. Der ist viel. Und kurz. Oder auch l?nger. Und kam Teilweise oben schon mal vor. So wie dieses hier: 'Im folgenden Kapitel sollen die physikalischen Grundlagen erl?utert werden, auf denen die Arbeit basiert'.  Und dann geht es wieder von norne los. Hier kommt jetzt massenhaft Text. Der ist viel. Und kurz. Oder auch l?nger. Und kam Teilweise oben schon mal vor. So wie dieses hier: 'Im folgenden Kapitel sollen die physikalischen Grundlagen erl?utert werden, auf denen die Arbeit basiert'.  Und dann geht es wieder von norne los. Hier kommt jetzt massenhaft Text. Der ist viel. Und kurz. Oder auch l?nger. Und kam Teilweise oben schon mal vor. So wie dieses hier: 'Im folgenden Kapitel sollen die physikalischen Grundlagen erl?utert werden, auf denen die Arbeit basiert'.  Und dann geht es wieder von norne los. Hier kommt jetzt massenhaft Text. Der ist viel. Und kurz. Oder auch l?nger. Und kam Teilweise oben schon mal vor. So wie dieses hier: 'Im folgenden Kapitel sollen die physikalischen Grundlagen erl?utert werden, auf denen die Arbeit basiert'.  Und dann geht es wieder von norne los. Hier kommt jetzt massenhaft Text. Der ist viel. Und kurz. Oder auch l?nger. Und kam Teilweise oben schon mal vor. So wie dieses hier: 'Im folgenden Kapitel sollen die physikalischen Grundlagen erl?utert werden, auf denen die Arbeit basiert'.  Und dann geht es wieder von norne los. Hier kommt jetzt massenhaft Text. Der ist viel. Und kurz. Oder auch l?nger. Und kam Teilweise oben schon mal vor. So wie dieses hier: 'Im folgenden Kapitel sollen die physikalischen Grundlagen erl?utert werden, auf denen die Arbeit basiert'.  Und dann geht es wieder von norne los. Hier kommt jetzt massenhaft Text. Der ist viel. Und kurz. Oder auch l?nger. Und kam Teilweise oben schon mal vor. So wie dieses hier: 'Im folgenden Kapitel sollen die physikalischen Grundlagen erl?utert werden, auf denen die Arbeit basiert'.  Und dann geht es wieder von norne los. Hier kommt jetzt massenhaft Text. Der ist viel. Und kurz. Oder auch l?nger. Und kam Teilweise oben schon mal vor. So wie dieses hier: 'Im folgenden Kapitel sollen die physikalischen Grundlagen erl?utert werden, auf denen die Arbeit basiert'.  Und dann geht es wieder von norne los. Hier kommt jetzt massenhaft Text. Der ist viel. Und kurz. Oder auch l?nger. Und kam Teilweise oben schon mal vor. So wie dieses hier: 'Im folgenden Kapitel sollen die physikalischen Grundlagen erl?utert werden, auf denen die Arbeit basiert'.  Und dann geht es wieder von norne los. Hier kommt jetzt massenhaft Text. Der ist viel. Und kurz. Oder auch l?nger. Und kam Teilweise oben schon mal vor. So wie dieses hier: 'Im folgenden Kapitel sollen die physikalischen Grundlagen erl?utert werden, auf denen die Arbeit basiert'.  Und dann geht es wieder von norne los. Hier kommt jetzt massenhaft Text. Der ist viel. Und kurz. Oder auch l?nger. Und kam Teilweise oben schon mal vor. So wie dieses hier: 'Im folgenden Kapitel sollen die physikalischen Grundlagen erl?utert werden, auf denen die Arbeit basiert'.  Und dann geht es wieder von norne los. Hier kommt jetzt massenhaft Text. Der ist viel. Und kurz. Oder auch l?nger. Und kam Teilweise oben schon mal vor. So wie dieses hier: 'Im folgenden Kapitel sollen die physikalischen Grundlagen erl?utert werden, auf denen die Arbeit basiert'.  Und dann geht es wieder von norne los. Hier kommt jetzt massenhaft Text. Der ist viel. Und kurz. Oder auch l?nger. Und kam Teilweise oben schon mal vor. So wie dieses hier: 'Im folgenden Kapitel sollen die physikalischen Grundlagen erl?utert werden, auf denen die Arbeit basiert'.  Und dann geht es wieder von norne los. Hier kommt jetzt massenhaft Text. Der ist viel. Und kurz. Oder auch l?nger. Und kam Teilweise oben schon mal vor. So wie dieses hier: 'Im folgenden Kapitel sollen die physikalischen Grundlagen erl?utert werden, auf denen die Arbeit basiert'.  Und dann geht es wieder von norne los. Hier kommt jetzt massenhaft Text. Der ist viel. Und kurz. Oder auch l?nger. Und kam Teilweise oben schon mal vor. So wie dieses hier: 'Im folgenden Kapitel sollen die physikalischen Grundlagen erl?utert werden, auf denen die Arbeit basiert'.  Und dann geht es wieder von norne los. 


Hier kommt jetzt massenhaft Text. Der ist viel. Und kurz. Oder auch l?nger. Und kam Teilweise oben schon mal vor. So wie dieses hier: 'Im folgenden Kapitel sollen die physikalischen Grundlagen erl?utert werden, auf denen die Arbeit basiert'.  Und dann geht es wieder von norne los. Hier kommt jetzt massenhaft Text. Der ist viel. Und kurz. Oder auch l?nger. Und kam Teilweise oben schon mal vor. So wie dieses hier: 'Im folgenden Kapitel sollen die physikalischen Grundlagen erl?utert werden, auf denen die Arbeit basiert'.  Und dann geht es wieder von norne los. Hier kommt jetzt massenhaft Text. Der ist viel. Und kurz. Oder auch l?nger. Und kam Teilweise oben schon mal vor. So wie dieses hier: 'Im folgenden Kapitel sollen die physikalischen Grundlagen erl?utert werden, auf denen die Arbeit basiert'.  Und dann geht es wieder von norne los. Hier kommt jetzt massenhaft Text. Der ist viel. Und kurz. Oder auch l?nger. Und kam Teilweise oben schon mal vor. So wie dieses hier: 'Im folgenden Kapitel sollen die physikalischen Grundlagen erl?utert werden, auf denen die Arbeit basiert'.  Und dann geht es wieder von norne los. Hier kommt jetzt massenhaft Text. Der ist viel. Und kurz. Oder auch l?nger. Und kam Teilweise oben schon mal vor. So wie dieses hier: 'Im folgenden Kapitel sollen die physikalischen Grundlagen erl?utert werden, auf denen die Arbeit basiert'.  Und dann geht es wieder von norne los. Hier kommt jetzt massenhaft Text. Der ist viel. Und kurz. Oder auch l?nger. Und kam Teilweise oben schon mal vor. So wie dieses hier: 'Im folgenden Kapitel sollen die physikalischen Grundlagen erl?utert werden, auf denen die Arbeit basiert'.  Und dann geht es wieder von norne los. Hier kommt jetzt massenhaft Text. Der ist viel. Und kurz. Oder auch l?nger. Und kam Teilweise oben schon mal vor. So wie dieses hier: 'Im folgenden Kapitel sollen die physikalischen Grundlagen erl?utert werden, auf denen die Arbeit basiert'.  Und dann geht es wieder von norne los. Hier kommt jetzt massenhaft Text. Der ist viel. Und kurz. Oder auch l?nger. Und kam Teilweise oben schon mal vor. So wie dieses hier: 'Im folgenden Kapitel sollen die physikalischen Grundlagen erl?utert werden, auf denen die Arbeit basiert'.  Und dann geht es wieder von norne los. Hier kommt jetzt massenhaft Text. Der ist viel. Und kurz. Oder auch l?nger. Und kam Teilweise oben schon mal vor. So wie dieses hier: 'Im folgenden Kapitel sollen die physikalischen Grundlagen erl?utert werden, auf denen die Arbeit basiert'.  Und dann geht es wieder von norne los. Hier kommt jetzt massenhaft Text. Der ist viel. Und kurz. Oder auch l?nger. Und kam Teilweise oben schon mal vor. So wie dieses hier: 'Im folgenden Kapitel sollen die physikalischen Grundlagen erl?utert werden, auf denen die Arbeit basiert'.  Und dann geht es wieder von norne los. Hier kommt jetzt massenhaft Text. Der ist viel. Und kurz. Oder auch l?nger. Und kam Teilweise oben schon mal vor. So wie dieses hier: 'Im folgenden Kapitel sollen die physikalischen Grundlagen erl?utert werden, auf denen die Arbeit basiert'.  Und dann geht es wieder von norne los. Hier kommt jetzt massenhaft Text. Der ist viel. Und kurz. Oder auch l?nger. Und kam Teilweise oben schon mal vor. So wie dieses hier: 'Im folgenden Kapitel sollen die physikalischen Grundlagen erl?utert werden, auf denen die Arbeit basiert'.  Und dann geht es wieder von norne los. Hier kommt jetzt massenhaft Text. Der ist viel. Und kurz. Oder auch l?nger. Und kam Teilweise oben schon mal vor. So wie dieses hier: 'Im folgenden Kapitel sollen die physikalischen Grundlagen erl?utert werden, auf denen die Arbeit basiert'.  Und dann geht es wieder von norne los. Hier kommt jetzt massenhaft Text. Der ist viel. Und kurz. Oder auch l?nger. Und kam Teilweise oben schon mal vor. So wie dieses hier: 'Im folgenden Kapitel sollen die physikalischen Grundlagen erl?utert werden, auf denen die Arbeit basiert'.  Und dann geht es wieder von norne los. Hier kommt jetzt massenhaft Text. Der ist viel. Und kurz. Oder auch l?nger. Und kam Teilweise oben schon mal vor. So wie dieses hier: 'Im folgenden Kapitel sollen die physikalischen Grundlagen erl?utert werden, auf denen die Arbeit basiert'.  Und dann geht es wieder von norne los. 


Hier kommt jetzt massenhaft Text. Der ist viel. Und kurz. Oder auch l?nger. Und kam Teilweise oben schon mal vor. So wie dieses hier: 'Im folgenden Kapitel sollen die physikalischen Grundlagen erl?utert werden, auf denen die Arbeit basiert'.  Und dann geht es wieder von norne los. Hier kommt jetzt massenhaft Text. Der ist viel. Und kurz. Oder auch l?nger. Und kam Teilweise oben schon mal vor. So wie dieses hier: 'Im folgenden Kapitel sollen die physikalischen Grundlagen erl?utert werden, auf denen die Arbeit basiert'.  Und dann geht es wieder von norne los. Hier kommt jetzt massenhaft Text. Der ist viel. Und kurz. Oder auch l?nger. Und kam Teilweise oben schon mal vor. So wie dieses hier: 'Im folgenden Kapitel sollen die physikalischen Grundlagen erl?utert werden, auf denen die Arbeit basiert'.  Und dann geht es wieder von norne los. Hier kommt jetzt massenhaft Text. Der ist viel. Und kurz. Oder auch l?nger. Und kam Teilweise oben schon mal vor. So wie dieses hier: 'Im folgenden Kapitel sollen die physikalischen Grundlagen erl?utert werden, auf denen die Arbeit basiert'.  Und dann geht es wieder von norne los. Hier kommt jetzt massenhaft Text. Der ist viel. Und kurz. Oder auch l?nger. Und kam Teilweise oben schon mal vor. So wie dieses hier: 'Im folgenden Kapitel sollen die physikalischen Grundlagen erl?utert werden, auf denen die Arbeit basiert'.  Und dann geht es wieder von norne los. Hier kommt jetzt massenhaft Text. Der ist viel. Und kurz. Oder auch l?nger. Und kam Teilweise oben schon mal vor. So wie dieses hier: 'Im folgenden Kapitel sollen die physikalischen Grundlagen erl?utert werden, auf denen die Arbeit basiert'.  Und dann geht es wieder von norne los. Hier kommt jetzt massenhaft Text. Der ist viel. Und kurz. Oder auch l?nger. Und kam Teilweise oben schon mal vor. So wie dieses hier: 'Im folgenden Kapitel sollen die physikalischen Grundlagen erl?utert werden, auf denen die Arbeit basiert'.  Und dann geht es wieder von norne los. Hier kommt jetzt massenhaft Text. Der ist viel. Und kurz. Oder auch l?nger. Und kam Teilweise oben schon mal vor. So wie dieses hier: 'Im folgenden Kapitel sollen die physikalischen Grundlagen erl?utert werden, auf denen die Arbeit basiert'.  Und dann geht es wieder von norne los. Hier kommt jetzt massenhaft Text. Der ist viel. Und kurz. Oder auch l?nger. Und kam Teilweise oben schon mal vor. So wie dieses hier: 'Im folgenden Kapitel sollen die physikalischen Grundlagen erl?utert werden, auf denen die Arbeit basiert'.  Und dann geht es wieder von norne los. Hier kommt jetzt massenhaft Text. Der ist viel. Und kurz. Oder auch l?nger. Und kam Teilweise oben schon mal vor. So wie dieses hier: 'Im folgenden Kapitel sollen die physikalischen Grundlagen erl?utert werden, auf denen die Arbeit basiert'.  Und dann geht es wieder von norne los. Hier kommt jetzt massenhaft Text. Der ist viel. Und kurz. Oder auch l?nger. Und kam Teilweise oben schon mal vor. So wie dieses hier: 'Im folgenden Kapitel sollen die physikalischen Grundlagen erl?utert werden, auf denen die Arbeit basiert'.  Und dann geht es wieder von norne los. 
